\documentclass{crypto-exercise}
\author{Sven Laur}
\tags{Computational Diffie-Hellman problem, random self-reducibility, majority-voting}

\begin{document}
\begin{exercise}{Majority voting of CDH}
  Let $\GG$ be a finite group such that all elements $y\in\GG$ can be
  expressed as powers of $g\in\GG$. Then the Computational
  Diffie-Hellman (CDH) problem is following. Given $x=g^a$ and $y=g^b$, find a group element $z=g^{ab}$. 
  Assume that there exists an algorithm $\AD$ that, on input $(g^x, g^y)\in \mathbb{G}^2$, outputs $g^{xy}$ with probability $\varepsilon$, and outputs a uniformly random element of $\mathbb{G}\setminus\{g^{xy}\}$ otherwise. Such an algorithm can be obtained by converting an arbitrary CDH solver into a randomly self-reduced CDH solver.
Consider now the following CDH amplification strategy $\ADB$:
\begin{itemize}
  \item Run $\AD$ a total of $\ell$ times to obtain list of outputs $w_1, \ldots, w_\ell$.
  \item Output $\bot$ if all candidate outputs are distinct.
  \item Output $\bot$ if more than one group element appears at least twice among the candidates.
  \item Otherwise, output the unique value $w$ that occurs at least twice in the list of candidates.
\end{itemize}
Estimate what is the probability that the answer is correct.   How does this probability scale wrt $\ell$.
\end{exercise}
\begin{solution}
	
Let us define two non-disjoint events $\mathsf{CDHPair}$ and $\mathsf{RandPair}$ where the first denotes the event that $g^{xy}$ occures more than once in the list and the second that a random $w\neq g^{xy}$ occurs in the list. Then we can estimate 
\begin{align*}
\pr{\ADB(g^x,g^y)=g^{xy}}\geq \pr{\mathsf{CDHPair}}-\pr{\mathsf{RandPair}}\enspace
\end{align*}
and we can use inclusion-exclusion principle to lower bound the first and upper bound the second.	
\textbf{Hint:} If this gives overly conservative estimates, then you can use precise expression 
\begin{align*}
\pr{\ADB(g^x,g^y)=g^{xy}}= \pr{\mathsf{CDHPair}}\cdot\pr{\neg\mathsf{RandPair}|\mathsf{CDHPair}}\enspace
\end{align*}
and lowerbound the second term.

\end{solution}
\end{document}
